%% Appendix D -- supplementary budget analyses (moved from Sects. 4.4 and 5.1 on 2026-09-29; numbers unchanged).
\section{Supplementary budget analyses}
\label{app:D}

\subsection{Choice of denominator}
\label{app:D1}

The ratio of total to reported uncertainty in Sect.~\ref{sec:4.4} uses the per-scan fit-propagated $\sigma$. Three
other quantities are available, and none of them answers the same question (Table~\ref{tab:denominators}).

\begin{table}[t]
\caption{Candidate denominators for the ratio of total to reported $\Sigma$ANs uncertainty: value (ppb), time resolution, and the ratio of the total uncertainty (Table~\ref{tab:2}, robust scale to standard deviation) to each. The zero-air value is from Sect.~\ref{sec:7.2}.}
\label{tab:denominators}
\begin{tabular}{llll}
\tophline
denominator & ppb & resolution & total / denominator \\
\middlehline
fit-propagated $\sigma$, per scan (used) & \fieldnum{0.0541} & 60\,s = 1 scan & \fieldnum{1.63}--\fieldnum{2.50} \\
uncertainty column shipped with the product & \fieldnum{0.0806} & 5\,min, systematic & \fieldnum{1.32}--\fieldnum{1.83} \\
fit-propagated $\sigma$, 5\,min bin & \fieldnum{0.025} & 5\,min & not comparable \\
zero-air scatter, SD & \fieldnum{0.0856} & 34-scan block means & not comparable \\
\bottomhline
\end{tabular}
\end{table}

The column shipped with the product is a gain systematic rather than a measurement uncertainty. Regressed against
the channel NO$_2$ amount over \fieldnum{14\,974} records it is $\fieldnum{0.04209}\,[\mathrm{NO_2}] + \fieldnum{0.00025}$\,ppb with
$r = \fieldnum{+0.9997}$, and the product header states that it covers day-to-day inter-channel gain only and must not be used
alone to assess detection significance; it does not bound the absolute scale either (Sect.~\ref{sec:3.6}). Comparing
the structural budget against it asks whether structure is comparable to the gain term, not whether it exceeds the
precision the retrieval claims. The five-minute row is the fit uncertainty averaged over five scans -- over the
deployment it is $\fieldnum{0.0558}/\sqrt{5}$ -- and dividing it into a structural scale measured at 60\,s would mix resolutions
in the opposite direction. The zero-air scatter is the precision floor, but it is measured on 34-scan block means,
each retrieved against a reference interpolated across the removed block, so it combines a different scan count with
reference-interpolation error and is not a single-scan quantity.

The per-scan fit uncertainty is the median, over \fieldnum{8832} of the budget records, of the NO$_2$ uncertainty written by
the operational fit of each channel propagated through the channel difference, $\sqrt{\sigma_{300}^2 +
g'^2\sigma_{180}^2}$. Over the whole deployment, with the $g'$ of each calibration state, the same median is
\fieldnum{0.0558}\,ppb (total \fieldnum{1.59}--\fieldnum{2.43} times), and at five-minute resolution the fit uncertainty lies between \fieldnum{0.0217}
and \fieldnum{0.0326}\,ppb (5th and 95th percentiles).

\subsection{Coupling of the structural terms to the wavelength shift}
\label{app:D2}

The structural terms of Sect.~\ref{sec:4} differ in how strongly they couple to the shift
(Table~\ref{tab:shift_degeneracy}). An error in the reflectivity changes the magnitude of the extinction and little
else: it moves the shift by a median of \fieldnum{0.006}\,px (90th percentile \fieldnum{0.027}\,px), although the largest displacement
reaches \fieldnum{0.285}\,px. An error in the zero-air reference or in the etalon frequency changes the shape of the extinction
and is partly absorbed by the wavelength axis. The shape statistic of Sect.~\ref{sec:6.1a} gives the same distinction
independently: the reflectivity is a pure amplitude term, and the other two compete with the nonlinear parameters.

\begin{table}[t]
\caption{Change in the fitted wavelength shift (px) when each structural term of Sect.~\ref{sec:4} is perturbed alone, and
when all three are perturbed together: median, 90th percentile and maximum over the \fieldnum{8847} clock-aligned budget
records of Table~\ref{tab:2} (for each $\Sigma$ANs record, the larger of the two channels' changes).}
\label{tab:shift_degeneracy}
\begin{tabular}{llll}
\tophline
perturbation & median & $p_{90}$ & max \\
\middlehline
zero-air reference & \fieldnum{0.187}\,px & \fieldnum{0.826} & \fieldnum{4.388} \\
reflectivity & \fieldnum{0.006} & \fieldnum{0.027} & \fieldnum{0.285} \\
etalon frequency & \fieldnum{0.154} & \fieldnum{0.347} & \fieldnum{2.941} \\
all three, joint & \fieldnum{0.243} & \fieldnum{0.766} & \fieldnum{4.276} \\
\bottomhline
\end{tabular}
\end{table}

The same propagation shows that the budget of Sect.~\ref{sec:4} is not an artefact of basin hopping in the shift. Records whose shift moved by more than a pixel
under perturbation are \fieldnum{6.2}\,\% of the total and carry \fieldnum{5.0}\,\% of the variance, and removing them raises the budget
(rSD \fieldnum{0.0693} to \fieldnum{0.0706}\,ppb, SD \fieldnum{0.1236} to \fieldnum{0.1244}\,ppb). These are the low-amplification records, where the
absorption is deep and the retrieval well conditioned.
